% ============================================================
%  Chapter 3 -- Addition and Multiplication of Series,
%               Rearrangements: 3.47-3.55
% ============================================================
\rsection{Addition and Multiplication of Series}

\begin{signpost}[Which finite-sum rules survive the limit]
The chapter's closing question: which of the ordinary rules of arithmetic
survive when ``sum'' means ``limit of partial sums''? Addition survives
unconditionally (3.47). Multiplication survives if at least one factor
converges absolutely (Mertens, 3.50), can fail for two conditionally
convergent factors (3.49), yet can never converge to the \emph{wrong} value
(Abel, 3.51). Rearrangement --- the commutative law --- is the most fragile of
all, and the next section is devoted to its spectacular failure.
\end{signpost}

\ritem{3.47}{Theorem}
\emph{If $\Sigma a_n = A$ and $\Sigma b_n = B$, then
$\Sigma(a_n+b_n) = A+B$, and $\Sigma ca_n = cA$, for any fixed $c$.}

\begin{proof}
Let
\[ A_n = \sum_{k=0}^{n} a_k, \qquad B_n = \sum_{k=0}^{n} b_k. \]
Then
\[ A_n + B_n = \sum_{k=0}^{n} (a_k+b_k). \]
Since $\lim_{n\to\infty} A_n = A$ and $\lim_{n\to\infty} B_n = B$, we see that
\[ \lim_{n\to\infty} (A_n+B_n) = A+B. \]
The proof of the second assertion is even simpler.
\end{proof}

Thus two convergent series may be added term by term, and the resulting series
converges to the sum of the two series. The situation becomes more complicated
when we consider multiplication of two series. To begin with, we have to define
the product. This can be done in several ways; we shall consider the so-called
``Cauchy product.''

\ritem{3.48}{Definition}
Given $\Sigma a_n$ and $\Sigma b_n$, we put
\[ c_n = \sum_{k=0}^{n} a_kb_{n-k} \qquad (n = 0,1,2,\dots) \]
and call $\Sigma c_n$ the \emph{product} of the two given series.

This definition may be motivated as follows. If we take two power series
$\Sigma a_nz^n$ and $\Sigma b_nz^n$, multiply them term by term, and collect
terms containing the same power of $z$, we get
\begin{align*}
  \sum_{n=0}^{\infty} a_nz^n \cdot \sum_{n=0}^{\infty} b_nz^n
    &= (a_0 + a_1z + a_2z^2 + \cdots)(b_0 + b_1z + b_2z^2 + \cdots)\\
    &= a_0b_0 + (a_0b_1 + a_1b_0)z\\
    &\qquad + (a_0b_2 + a_1b_1 + a_2b_0)z^2 + \cdots\\
    &= c_0 + c_1z + c_2z^2 + \cdots.
\end{align*}
Setting $z = 1$, we arrive at the above definition.\kn{The $c_n$ collect the
products $a_kb_j$ along the anti-diagonals $k+j = n$ of the multiplication
array --- the same anti-diagonals as in the countability proof 2.12, and for
the same reason: each one is finite. The definition is rigged for power
series, and 3.51's proof will return there.}

\begin{center}
\begin{tikzpicture}[x=1mm, y=1mm]
  \fill[black!8] (0,0) -- (30,0) -- (0,30) -- cycle;
  \fill[pattern=north east lines, pattern color=black!40]
    (30,0) -- (30,30) -- (0,30) -- cycle;
  \draw[guide] (0,0) rectangle (30,30);
  \draw[thin] (30,0) -- (0,30);
  \draw[thin] (-1.5,19.5) -- (19.5,-1.5);
  \foreach \i in {0,6,...,30}{\foreach \j in {0,6,...,30}{
    \node[ptfill, minimum size=1.7pt] at (\i,\j) {};}}
  \node[mlbl, anchor=north west] at (20,-1.8) {$c_3$};
  \node[mlbl, fill=white, inner sep=1pt] at (6.5,3.5) {$C_N$};
  \node[mlbl, fill=white, inner sep=1pt] at (21,26.5) {$A_NB_N - C_N$};
  \node[mlbl, anchor=north] at (30,-2.4) {$i=N$};
  \node[mlbl, anchor=east] at (-2.4,30) {$j=N$};
\end{tikzpicture}
\end{center}
\figcap{The multiplication array $a_ib_j$, one dot per product. Each $c_n$
strings an anti-diagonal (shown: $i+j=3$), and every anti-diagonal is finite
--- the definition's whole content. The shaded triangle $i+j \le N$ sums to
$C_N$; the dashed square sums to $A_NB_N$. Everything about $\Sigma c_n$
versus $AB$ is the fate of the hatched difference: Mertens (3.50) shows
absolute convergence of one factor kills it, while in Example 3.49 it carries
mass ${\approx}2$ forever. These are the anti-diagonals of 2.12's walk,
finite for the same reason.}

\ritem{3.49}{Example}
If
\[ A_n = \sum_{k=0}^{n} a_k, \qquad B_n = \sum_{k=0}^{n} b_k, \qquad
   C_n = \sum_{k=0}^{n} c_k, \]
and $A_n \to A$, $B_n \to B$, then it is not at all clear that $(C_n)$ will
converge to $AB$, since we do not have $C_n = A_nB_n$. The dependence of
$(C_n)$ on $(A_n)$ and $(B_n)$ is quite a complicated one (see the proof of
Theorem 3.50). We shall now show that the product of two convergent series may
actually diverge.

The series
\[ \sum_{n=0}^{\infty} \frac{(-1)^n}{\sqrt{n+1}}
   = 1 - \frac{1}{\sqrt2} + \frac{1}{\sqrt3} - \frac{1}{\sqrt4} + \cdots \]
converges (Theorem 3.43). We form the product of this series with itself and
obtain
\begin{align*}
  \sum_{n=0}^{\infty} c_n = 1
    &- \left(\frac{1}{\sqrt2}+\frac{1}{\sqrt2}\right)
     + \left(\frac{1}{\sqrt3}+\frac{1}{\sqrt2\sqrt2}+\frac{1}{\sqrt3}\right)\\
    &- \left(\frac{1}{\sqrt4}+\frac{1}{\sqrt3\sqrt2}+\frac{1}{\sqrt2\sqrt3}
     +\frac{1}{\sqrt4}\right) + \cdots,
\end{align*}
so that
\[ c_n = (-1)^n \sum_{k=0}^{n} \frac{1}{\sqrt{(n-k+1)(k+1)}}. \]
Since
\[ (n-k+1)(k+1) = \left(\frac n2+1\right)^2 - \left(\frac n2-k\right)^2
   \le \left(\frac n2+1\right)^2, \]
we have
\[ |c_n| \ge \sum_{k=0}^{n} \frac{2}{n+2} = \frac{2(n+1)}{n+2}, \]
so that the condition $c_n \to 0$, which is necessary for the convergence of
$\Sigma c_n$, is not satisfied.\kn{Why it fails: each $c_n$ averages $n+1$
products of terms of size about $1/\sqrt n$ each, so the products are of size
about $1/n$ and their sum stays near $2$. The factor-series converge only by
cancellation, and the Cauchy product lines up terms so that the cancellation
is destroyed. Convergence-by-cancellation is not robust under
rearrangement-like operations --- the theme of the whole final section.}

In view of the next theorem, due to Mertens, we note that we have here
considered the product of two conditionally convergent series.

\ritem{3.50}{Theorem}
\emph{Suppose}
\begin{enumerate}[label=(\alph*), leftmargin=2.2em, itemsep=1pt, topsep=3pt]
  \item \emph{$\sum_{n=0}^{\infty} a_n$ converges absolutely,}
  \item \emph{$\sum_{n=0}^{\infty} a_n = A$,}
  \item \emph{$\sum_{n=0}^{\infty} b_n = B$,}
  \item \emph{$c_n = \sum_{k=0}^{n} a_kb_{n-k}$ \quad $(n = 0,1,2,\dots)$.}
\end{enumerate}
\emph{Then}
\[ \sum_{n=0}^{\infty} c_n = AB. \]

That is, the product of two convergent series converges, and to the right
value, if at least one of the two series converges absolutely.

\begin{proof}
Put
\[ A_n = \sum_{k=0}^{n} a_k, \qquad B_n = \sum_{k=0}^{n} b_k, \qquad
   C_n = \sum_{k=0}^{n} c_k, \qquad \beta_n = B_n - B. \]
Then
\begin{align*}
  C_n &= a_0b_0 + (a_0b_1+a_1b_0) + \cdots + (a_0b_n+a_1b_{n-1}+\cdots+a_nb_0)\\
      &= a_0B_n + a_1B_{n-1} + \cdots + a_nB_0\\
      &= a_0(B+\beta_n) + a_1(B+\beta_{n-1}) + \cdots + a_n(B+\beta_0)\\
      &= A_nB + a_0\beta_n + a_1\beta_{n-1} + \cdots + a_n\beta_0.
\end{align*}\kn{The first rewriting is the anti-diagonal picture again: summing
$c_0,\dots,c_n$ sweeps the triangle $k+j \le n$ of the array, which can be
sliced as $a_k$ times the partial sum $B_{n-k}$. Then each $B_{n-k}$ is split
into limit plus error, $B + \beta_{n-k}$, isolating the main term $A_nB$. The
proof will live or die on the error convolution $\gamma_n$.}
Put
\[ \gamma_n = a_0\beta_n + a_1\beta_{n-1} + \cdots + a_n\beta_0. \]
We wish to show that $C_n \to AB$. Since $A_nB \to AB$, it suffices to show
that
\begin{equation}
  \lim_{n\to\infty} \gamma_n = 0. \tag{3.21}
\end{equation}
Put
\[ \alpha = \sum_{n=0}^{\infty} |a_n|. \]
[It is here that we use (a).] Let $\varepsilon > 0$ be given. By (c),
$\beta_n \to 0$. Hence we can choose $N$ such that $|\beta_n| \le \varepsilon$
for $n \ge N$, in which case
\begin{align*}
  |\gamma_n| &\le |\beta_0a_n + \cdots + \beta_Na_{n-N}|
              + |\beta_{N+1}a_{n-N-1} + \cdots + \beta_na_0|\\
             &\le |\beta_0a_n + \cdots + \beta_Na_{n-N}| + \varepsilon\alpha.
\end{align*}
Keeping $N$ fixed and letting $n \to \infty$, we get
\[ \limsup_{n\to\infty} |\gamma_n| \le \varepsilon\alpha, \]
since $a_k \to 0$ as $k \to \infty$. Since $\varepsilon$ is arbitrary, (3.21)
follows.\kn{The split-the-sum device, and the same two-parameter discipline as
in 3.31: the \emph{old} errors $\beta_0,\dots,\beta_N$ are finitely many and
ride multiples of $a_n \to 0$; the \emph{recent} errors are small and ride the
absolutely convergent tail, total at most $\varepsilon\alpha$. Freeze $N$, let
$n$ run, then shrink $\varepsilon$. Without (a) the constant $\alpha$ does not
exist, and 3.49 shows the theorem genuinely fails.}
\end{proof}

\begin{center}
\begin{tikzpicture}[x=1mm, y=1mm]
  \node[rmbox, text width=30mm] (m0) at (0,24)
    {slice anti-diagonals:\\$C_n = A_nB + \gamma_n$};
  \node[rmbox, text width=28mm] (m1) at (60,24)
    {goal shrinks to\\$\gamma_n \to 0$};
  \draw[rmarrow] (m0) -- (m1);
  \node[rmbox, text width=31mm] (m2) at (0,4)
    {head: $\beta_0,\dots,\beta_N$\\ride $a_n \to 0$};
  \node[rmbox, text width=31mm] (m3) at (60,4)
    {tail: $|{\cdot}| \le \varepsilon\alpha$\\ \textit{(absolute convergence)}};
  \draw[rmarrow] (m1) -- (m2);
  \draw[rmarrow] (m1) -- (m3);
  \node[rmbox, text width=34mm] (m4) at (30,-15)
    {$\limsup|\gamma_n| \le \varepsilon\alpha$, all $\varepsilon$:\\
     $C_n \to AB$};
  \draw[rmarrow] (m2) -- (m4);
  \draw[rmarrow] (m3) -- (m4);
\end{tikzpicture}
\end{center}
\figcap{Mertens as a flow diagram: isolate the error convolution, split it at
$N$, and let each half die for its own reason. Hypothesis (a) funds exactly one
box --- the tail.}

Another question which may be asked is whether the series $\Sigma c_n$, if
convergent, must have the sum $AB$. Abel showed that the answer is in the
affirmative.

\ritem{3.51}{Theorem}
\emph{If the series $\Sigma a_n$, $\Sigma b_n$, $\Sigma c_n$ converge to $A$,
$B$, $C$, and $c_n = a_0b_n + \cdots + a_nb_0$, then $C = AB$.}

Here no assumption is made concerning absolute convergence. We shall give a
simple proof (which depends on the continuity of power series) after Theorem
8.2.\kd{A promissory note worth remembering when it is paid. The future proof:
the three power series $\Sigma a_nx^n$, $\Sigma b_nx^n$, $\Sigma c_nx^n$
converge absolutely for $0 < x < 1$, where Mertens applies and gives
$f(x)g(x) = h(x)$; then let $x \to 1$ and invoke Abel's continuity theorem
(8.2) three times. The pattern --- prove an identity where convergence is
absolute, then reach the boundary by continuity --- is a first taste of a
standard analytic method.}

\rsection{Rearrangements}

\ritem{3.52}{Definition}
Let $(k_n)$, $n = 1,2,3,\dots$, be a sequence in which every positive integer
appears once and only once (that is, $(k_n)$ is a 1-1 function from $\N$ onto
$\N$, in the notation of Definition 2.2). Putting
\[ a_n' = a_{k_n} \qquad (n = 1,2,3,\dots), \]
we say that $\Sigma a_n'$ is a \emph{rearrangement} of $\Sigma a_n$.

If $(s_n)$, $(s_n')$ are the sequences of partial sums of $\Sigma a_n$,
$\Sigma a_n'$, it is easily seen that, in general, these two sequences consist
of entirely different numbers. We are thus led to the problem of determining
under what conditions all rearrangements of a convergent series will converge,
and whether the sums are necessarily the same.

\ritem{3.53}{Example}
Consider the convergent series
\begin{equation}
  1 - \frac12 + \frac13 - \frac14 + \frac15 - \frac16 + \cdots \tag{3.22}
\end{equation}
and one of its rearrangements
\begin{equation}
  1 + \frac13 - \frac12 + \frac15 + \frac17 - \frac14 + \frac19 + \frac{1}{11}
  - \frac16 + \cdots \tag{3.23}
\end{equation}
in which two positive terms are always followed by one negative. If $s$ is the
sum of (3.22), then
\[ s < 1 - \frac12 + \frac13 = \frac56. \]\kn{Why this bound: (3.22) is
alternating with decreasing terms, so by the zigzag picture at 3.43 the
partial sums oscillate around $s$, and every \emph{odd} partial sum
overestimates it --- $s < s_3 = 5/6$.}
Since
\[ \frac{1}{4k-3} + \frac{1}{4k-1} - \frac{1}{2k} > 0 \]
for $k \ge 1$, we see that $s_3' < s_6' < s_9' < \cdots$, where $s_n'$ is the
$n$th partial sum of (3.23).

Hence
\[ \limsup_{n\to\infty} s_n' > s_3' = \frac56, \]
so that (3.23) certainly does not converge to $s$ [we leave it to the reader to
verify that (3.23) does, however, converge].\kn{The same numbers, added in a
different order, give a different sum: taking positives twice as fast as
negatives borrows against the divergent positive half, and each three-term
block adds a surplus. This is the concrete disaster; 3.54 shows it is total.}

This example illustrates the following theorem, due to Riemann.

\ritem{3.54}{Theorem}
\emph{Let $\Sigma a_n$ be a series of real numbers which converges, but not
absolutely. Suppose}
\[ -\infty \le \alpha \le \beta \le +\infty. \]
\emph{Then there exists a rearrangement $\Sigma a_n'$ with partial sums
$s_n'$ such that}
\begin{equation}
  \liminf_{n\to\infty} s_n' = \alpha, \qquad
  \limsup_{n\to\infty} s_n' = \beta. \tag{3.24}
\end{equation}

\begin{proof}
Let
\[ p_n = \frac{|a_n|+a_n}{2}, \qquad q_n = \frac{|a_n|-a_n}{2}
   \qquad (n = 1,2,3,\dots). \]
Then $p_n - q_n = a_n$, $p_n + q_n = |a_n|$, $p_n \ge 0$, $q_n \ge 0$. The
series $\Sigma p_n$, $\Sigma q_n$ must both diverge, for if both were
convergent, then
\[ \Sigma(p_n+q_n) = \Sigma|a_n| \]
would converge, contrary to hypothesis. Since
\[ \sum_{n=1}^{N} a_n = \sum_{n=1}^{N} (p_n-q_n)
   = \sum_{n=1}^{N} p_n - \sum_{n=1}^{N} q_n, \]
divergence of $\Sigma p_n$ and convergence of $\Sigma q_n$ (or vice versa)
implies divergence of $\Sigma a_n$, again contrary to
hypothesis.\kn{The anatomy of conditional convergence, and the engine of the
whole proof: split the series into its positive part and its negative part,
and \emph{both halves diverge}. The finite sum $s$ is the difference of two
infinities kept in balance by the original ordering. That balance is the only
thing a rearrangement can destroy --- and it can destroy it completely,
because each half is an inexhaustible reservoir.}

Now let $P_1, P_2, P_3, \dots$ denote the nonnegative terms of $\Sigma a_n$,
in the order in which they occur, and let $Q_1, Q_2, Q_3, \dots$ be the
absolute values of the negative terms of $\Sigma a_n$, also in their original
order.

The series $\Sigma P_n$, $\Sigma Q_n$ differ from $\Sigma p_n$, $\Sigma q_n$
only by zero terms, and are therefore divergent.

We shall construct sequences $(m_n)$, $(k_n)$ such that the series
\begin{equation}
\begin{split}
  P_1 + \cdots + P_{m_1} - Q_1 - \cdots - Q_{k_1} + P_{m_1+1} + \cdots\\
  {}+ P_{m_2} - Q_{k_1+1} - \cdots - Q_{k_2} + \cdots,
\end{split}
\tag{3.25}
\end{equation}
which clearly is a rearrangement of $\Sigma a_n$, satisfies (3.24).

Choose real-valued sequences $(\alpha_n)$, $(\beta_n)$ such that
$\alpha_n \to \alpha$, $\beta_n \to \beta$, $\alpha_n < \beta_n$,
$\beta_1 > 0$.

Let $m_1$, $k_1$ be the smallest positive integers such that
\begin{align*}
  P_1 + \cdots + P_{m_1} &> \beta_1,\\
  P_1 + \cdots + P_{m_1} - Q_1 - \cdots - Q_{k_1} &< \alpha_1;
\end{align*}
let $m_2$, $k_2$ be the smallest positive integers such that
\begin{align*}
  P_1 + \cdots + P_{m_1} - Q_1 - \cdots - Q_{k_1} + P_{m_1+1} + \cdots
    + P_{m_2} &> \beta_2,\\
  P_1 + \cdots + P_{m_1} - Q_1 - \cdots - Q_{k_1} + P_{m_1+1} + \cdots
    + P_{m_2}\hspace{2.2em}&\\
  {}- Q_{k_1+1} - \cdots - Q_{k_2} &< \alpha_2,
\end{align*}
and continue in this way. This is possible since $\Sigma P_n$ and $\Sigma Q_n$
diverge.

If $x_n$, $y_n$ denote the partial sums of (3.25) whose last terms are
$P_{m_n}$, $-Q_{k_n}$, then
\[ |x_n - \beta_n| \le P_{m_n}, \qquad |y_n - \alpha_n| \le Q_{k_n}. \]
Since $P_n \to 0$ and $Q_n \to 0$ as $n \to \infty$, we see that
$x_n \to \beta$, $y_n \to \alpha$.\kn{The overshoot estimate is where
convergence of the \emph{original} series earns its keep: $a_n \to 0$ (3.23),
so the reservoir terms $P_n$, $Q_n$ shrink, and each crossing of the target
overshoots by at most the last term used. The steering wheel: pay out
positives until just past $\beta_n$, then negatives until just below
$\alpha_n$ --- ``smallest integer such that'' makes each overshoot minimal.}

Finally, it is clear that no number less than $\alpha$ or greater than
$\beta$ can be a subsequential limit of the partial sums (3.25).
\end{proof}

\begin{center}
\begin{tikzpicture}[x=1mm, y=1mm]
  \node[rmbox, text width=29mm] (r0) at (0,22)
    {conditional convergence:\\$\Sigma P_n$, $\Sigma Q_n$ diverge};
  \node[rmbox, text width=28mm] (r1) at (43,22)
    {steer: positives past $\beta_n$,\\negatives below $\alpha_n$};
  \node[rmbox, text width=26mm] (r2) at (85,22)
    {always possible:\\reservoirs infinite};
  \draw[rmarrow] (r0) -- (r1);
  \draw[rmarrow] (r2) -- (r1);
  \node[rmbox, text width=28mm] (r3) at (43,0)
    {overshoots $\le P_{m_n},Q_{k_n}$\\$\to 0$ by 3.23};
  \draw[rmarrow] (r1) -- (r3);
  \node[rmbox, text width=30mm] (r4) at (0,0)
    {$\limsup s_n' = \beta$,\\$\liminf s_n' = \alpha$};
  \draw[rmarrow] (r3) -- (r4);
\end{tikzpicture}
\end{center}
\figcap{Riemann's theorem as a flow diagram: divergence of both reservoirs
makes the steering possible, convergence of the original series ($a_n \to 0$)
makes it accurate. Each hypothesis funds exactly one arrow.}

\begin{center}
\begin{tikzpicture}[x=1mm, y=1mm]
  \draw[axis] (-2,0) -- (96,0);
  \node[lbl, anchor=north] at (94,-2) {terms used};
  % target band
  \draw[thin] (0,20) -- (92,20);
  \node[mlbl, anchor=west] at (92.5,20) {$\beta$};
  \draw[thin] (0,7) -- (92,7);
  \node[mlbl, anchor=west] at (92.5,7) {$\alpha$};
  % sawtooth: rise past beta, fall below alpha, overshoot shrinking
  \draw[thin] (0,0) -- (10,10) -- (22,24.5) -- (30,4.5)
              -- (46,23) -- (56,5.5) -- (72,21.8) -- (82,6.2) -- (92,17);
  \node[ptfill, minimum size=2.2pt] at (22,24.5) {};
  \node[ptfill, minimum size=2.2pt] at (30,4.5) {};
  \node[ptfill, minimum size=2.2pt] at (46,23) {};
  \node[lbl, anchor=south] at (22,26) {overshoots shrink like $P_{m_n}$};
  \node[lbl, anchor=north] at (56,3.5) {undershoots shrink like $Q_{k_n}$};
\end{tikzpicture}
\end{center}
\figcap{Riemann's steering: positives push the partial sums just past
$\beta$, negatives pull them just below $\alpha$, forever. Divergence of both
reservoirs keeps the steering possible; $a_n \to 0$ makes the overshoots die,
so the upper and lower limits are exactly $\beta$ and $\alpha$. Take
$\alpha=\beta$ to hit any prescribed sum.}

\ritem{3.55}{Theorem}
\emph{If $\Sigma a_n$ is a series of complex numbers which converges
absolutely, then every rearrangement of $\Sigma a_n$ converges, and they all
converge to the same sum.}

\begin{proof}
Let $\Sigma a_n'$ be a rearrangement, with partial sums $s_n'$. Given
$\varepsilon > 0$, there exists an integer $N$ such that $m \ge n \ge N$
implies
\begin{equation}
  \sum_{i=n}^{m} |a_i| \le \varepsilon. \tag{3.26}
\end{equation}
Now choose $p$ such that the integers $1,2,\dots,N$ are all contained in the
set $k_1,k_2,\dots,k_p$ (we use the notation of Definition 3.52). Then if
$n > p$, the numbers $a_1,\dots,a_N$ will cancel in the difference
$s_n - s_n'$, so that $|s_n - s_n'| \le \varepsilon$, by (3.26). Hence
$(s_n')$ converges to the same sum as $(s_n)$.\kn{Absolute convergence makes
the tail \emph{unconditionally} negligible: (3.26) bounds any finite batch of
far-out terms regardless of order or sign, so once both partial sums have
swallowed $a_1,\dots,a_N$, they differ only by scraps from the tail. Compare
3.54, where the tail's two halves were infinite reservoirs; here the tail is
just small. This is the precise sense of Rudin's slogan at 3.46 that
absolutely convergent series behave like finite sums.}
\end{proof}

\begin{recap}[Chapter 3 in one paragraph]
Convergence of a sequence is defined by the $\varepsilon$--$N$ contract (3.1),
and completeness makes it checkable from inside: Cauchy in $\R^k$ means
convergent (3.11), monotone-and-bounded means convergent (3.14), and every
real sequence has $\limsup$ and $\liminf$ as its extreme subsequential limits
(3.17). A series is the sequence of its partial sums (3.21), so those three
facts become the Cauchy criterion (3.22), the boundedness criterion for
nonnegative series (3.24), and the root test (3.33); every test is ultimately
a comparison (3.25) with the geometric series (3.26), except summation by
parts (3.41--3.44), which alone reaches conditional convergence. The number
$e$ arrives by series, fast enough for irrationality to fall out (3.30--3.32).
The chapter's final lesson is the boundary drawn at 3.46 and dramatised at
3.54--3.55: absolutely convergent series behave like finite sums under
products and rearrangements; conditionally convergent series have sums that
depend on the order of the terms --- so much so that any prescribed value, or
none, can be arranged.
\end{recap}
